% !TeX root = ./../main.tex

\section{Kwantowa korekcja błędów}\label{sec:lr:qec}
\agentinfo{
Sekcja napisana w zadaniu PHD-TS-31 (2026-11-27) na podstawie notatki \agentref{eksperymenty/przeglad-literatury/sekcja-3.2-korekcja.md}; wpisy pod znacznikiem LR-3.2 w biblio.bib.
Każda liczba ma w tym pliku komentarz ze słowem ŹRÓDŁO (klucz, miejsce w pracy, wersja); liczby sprawdzono ponownie 2026-11-27 w otwartych wersjach prac, kopie w \agentref{eksperymenty/przeglad-literatury/_robocze/PHD-TS-31/cache}.
Uzgodnienia: wniosek o koszcie korekcji podawać tylko w postaci zawężonej (platforma, kod, docelowy błąd, data); nie pisać, że korekcja jest niemożliwa w erze NISQ ani że wymaga tysięcy kubitów na kubit logiczny bez tych zastrzeżeń.
Pozycje dasu_computing_2026, lavasani_quantum_2026 i gidney_how_2025 to preprinty: nie nazywać ich publikacjami, a przed oddaniem pracy sprawdzić, czy się ukazały.
Wyliczenie liczby kubitów logicznych na procesorze ibm_fez oraz ilorazy kubitów fizycznych na kubit obliczenia są własne i tak są opisane w tekście.
Świadomie niecytowane prace ze wstępu artykułu ICCS: aaronson_improved_2004 (symulacja obwodów stabilizatorowych, nie kody) i dymarsky_quantum_2021 (bez związku z kosztem korekcji).
Po recenzji (PHD-TS-90, 2026-11-27): zdanie o progu poprawione wcześniej przez orkiestratora; ocena Preskilla z 2018 r. nie jest już przedstawiana jako podważona (w źródle dotyczy bliskiej przyszłości i narzutu, a wniosek sekcji jest z nią zgodny); narzut kodu powierzchniowego ujednolicono z podrozdziałem 2.3 (rząd tysiąca do dziesięciu tysięcy kubitów); współczynnik Lambda nazwano w tekście; przy 101 kubitach kodu o dystansie 7 dopisano cztery kubity do usuwania wycieku, żeby liczba zgadzała się z wyliczeniem 97 kubitów ze wzoru; cztery zdania podzielono, a wnioski własne oznaczono w tekście; dla zachowania objętości usunięto zdanie przejściowe kończące sekcję (odesłanie do sekcji 3.3 jest w sekcji 3.1).
Rok w zwrotach typu praca opublikowana w 2025 r. to rok numeru czasopisma, zgodny z wpisem w bazie (google_quantum_2025: w sieci od grudnia 2024 r.; bluvstein_fault-tolerant_2026: w sieci od listopada 2025 r.).
Do rozważenia (recenzja, ustalenie D17): dwa preprinty z 2026 r. o koszcie rozkładu liczb 2048-bitowych przy innych architekturach (arXiv:2602.11457, arXiv:2605.03951) są znane tylko z tytułów, niesprawdzone i niecytowane.
DO DECYZJI AUTORA: brzmienie wniosku końcowego sekcji.
Do sekcji 3.8: według notatki (punkt 5.6) w przeszukanym zakresie nie znaleziono pracy stosującej korekcję do kwantowych reprezentacji obrazów; zdania tego nie ma w tekście sekcji 3.2, bo zakres wyszukiwania ma podać sekcja 3.8.
}

Korekcja jest jedyną z~rodzin metod zdefiniowanych w~podrozdziale~\ref{subsec:tb:qec:taxonomy}, która usuwa błędy ze stanu kwantowego w~trakcie obliczeń; kody i~twierdzenie progowe opisano w~podrozdziale~\ref{subsec:tb:qec:codes}, a~wykład teorii zawiera przegląd~\cite{terhal_quantum_2015}.
Tutaj pytanie jest węższe: jaką skalę i~jakość mają realizacje eksperymentalne z~lat 2021--2026 i~ile kosztuje kubit logiczny na urządzeniu liczącym około stu kubitów.

Odpowiedzi nie dają prace algebraiczne, które konstruują kody stabilizatorowe o~coraz lepszych parametrach $[[n,k,d]]$~\cite{nguyen_quantum_2018,nguyen_novel_2020,lv_explicit_2020}, ani symulacje progu, wykonane także dla błędów koherentnych, czyli systematycznych obrotów~\cite{bravyi_correcting_2018}: o~koszcie decydują obwód pomiaru syndromu, dekoder i~szum konkretnego sprzętu (uwaga własna).
% ŹRÓDŁO: nguyen_quantum_2018, nguyen_novel_2020, lv_explicit_2020: tylko streszczenia (Crossref, OpenAlex; notatka sekcji 3.2), według których wynikiem prac są konstrukcje i parametry kodów; ostatnia część zdania jest uwagą własną.
% ŹRÓDŁO: bravyi_correcting_2018, streszczenie, arXiv:1710.02270v1 (symulacja kodu powierzchniowego, ,,first error threshold estimates for several toy models of coherent noise'').

Pierwsze realizacje powtarzanej korekcji w~kodach o~dystansie~3 pokazały, że cykl korekcji jest wykonalny, lecz nie daje jeszcze zysku.
Na dziesięciokubitowym komputerze jonowym kod $[[7,1,3]]$ z~dekoderem czasu rzeczywistego miał błąd logiczny około $2{,}7\%$ na cykl~\cite{ryan-anderson_realization_2021}, a~kod powierzchniowy na $17$ kubitach nadprzewodzących około $3\%$ na cykl po odrzuceniu przebiegów z~wykrytym wyciekiem~\cite{krinner2022realizing}; w~obu pracach autorzy wskazują, że zrównanie błędu logicznego z~fizycznym wymaga poprawy sprzętu.
% ŹRÓDŁO: ryan-anderson_realization_2021, sekcja II, tekst przy rys. 3 i tab. 3 (2.70(6)% per QEC cycle), streszczenie (ten qubit, pseudo-threshold), arXiv:2107.07505v1.
% ŹRÓDŁO: krinner2022realizing, streszczenie (17 physical qubits, 3% per cycle when rejecting runs with leakage) i dyskusja (break-even may be in reach for modest improvements), arXiv:2112.03708v1.
Najbliższa urządzeniom używanym w~tej pracy jest demonstracja na 27-kubitowym procesorze IBM o~topologii ciężkiego sześciokąta: kubit logiczny o~dystansie~3 zajął $23$ kubity, a~błąd logiczny na rundę wyniósł bez postselekcji $0{,}059(2)$ dla stanu $\ket{0_L}$ i~$0{,}113(5)$ dla stanu $\ket{+_L}$~\cite{sundaresan_demonstrating_2023}.
% ŹRÓDŁO: sundaresan_demonstrating_2023, podpis rys. 1 (23 qubits), sekcja wyników przy rys. 3a (ibm_peekskill, 27 qubit Falcon; 0.059(2) i 0.113(5) without post-selection; po odrzuceniu wycieku 0.041(1) i 0.088(3)), arXiv:2203.07205v2.

Do progu zbliżono się w~2023~r., a~zejście poniżej niego wykazano w~latach 2024--2025.
Na 72-kubitowym procesorze kod powierzchniowy o~dystansie~5, zajmujący $49$ kubitów, był już nieznacznie lepszy od kodów o~dystansie~3: $2{,}914\%$ wobec $3{,}028\%$ błędu na cykl~\cite{google_suppressing_2023}.
% ŹRÓDŁO: google_suppressing_2023, streszczenie (2.914% +/- 0.016% wobec 3.028% +/- 0.023%), sekcja I (72-qubit device, 49-qubit distance-5 code) i sekcja VII (do zejścia poniżej progu elementy muszą się poprawić o co najmniej 20%), arXiv:2207.06431v2.
W~pracy opublikowanej w~2025~r. błąd logiczny malał $\Lambda=2{,}14\pm0{,}02$ raza przy zwiększeniu dystansu o~2 ($\Lambda$: współczynnik tłumienia błędu).
Kod o~dystansie~7, zajmujący $101$ kubitów 105-kubitowego procesora (w~tym cztery do usuwania wycieku), miał błąd $0{,}143\%$ na cykl i~przechowywał stan $2{,}4\pm0{,}3$ raza dłużej niż najlepszy kubit fizyczny~\cite{google_quantum_2025}; kubit logiczny trwalszy od składników układu uzyskano wcześniej w~kodzie bozonowym~\cite{sivak_real-time_2023}.
% ŹRÓDŁO: google_quantum_2025, streszczenie (Lambda = 2.14 +/- 0.02; 101-qubit distance-7 code; 0.143% +/- 0.003% per cycle; 2.4 +/- 0.3), sekcja I (105-qubit processor) i tekst główny o pamięci w kodzie o dystansie 7 (49 data qubits, 48 measure qubits, and 4 additional leakage removal qubits; nazwa ,,error suppression factor'' dla Lambda), arXiv:2408.13687v1.
% ŹRÓDŁO: sivak_real-time_2023, streszczenie (coherence gain G = 2.27 +/- 0.07 względem najlepszego składnika) i wstęp (kod GKP w oscylatorze), arXiv:2211.09116v1 (2022).
Korekcja poniżej progu jest więc od lat 2024--2025 faktem eksperymentalnym, co nie podważa oceny z~2018~r., że urządzenia ery NISQ nie będą z~niej korzystać z~powodu narzutu~\cite{preskill_quantum_2018}.
% ŹRÓDŁO: preskill_quantum_2018, sekcja 6.4 (,,because of the high overhead cost of quantum error correction, NISQ devices will not make use of it in the near term''), arXiv:1801.00862v3.
Wyniki te dotyczą bowiem pamięci jednego kubitu logicznego zajmującej cały procesor, a~zysk jest kilkukrotny: według ekstrapolacji autorów błąd $10^{-6}$ wymagałby dystansu~27 i~$1457$ kubitów fizycznych na kubit logiczny~\cite{google_quantum_2025}.
% ŹRÓDŁO: google_quantum_2025, sekcja VI Outlook (distance-27 logical qubit using 1457 physical qubits; ,,resource intensive in practice''), arXiv:2408.13687v1.

Narzut rzędu $10^3$--$10^4$ kubitów fizycznych na kubit logiczny (podrozdział~\ref{subsec:tb:qec:codes}) nie jest przy tym własnością korekcji, lecz kodu powierzchniowego na siatce planarnej.
W~symulacjach kod $[[144,12,12]]$ o~rzadkiej macierzy kontroli parzystości (qLDPC) chroni $12$ kubitów logicznych na $288$ kubitach fizycznych przez blisko milion cykli przy błędzie fizycznym $0{,}1\%$, do czego kod powierzchniowy potrzebowałby blisko $3000$ kubitów~\cite{bravyi_high-threshold_2024}; wymaga jednak połączeń dalekiego zasięgu, których topologia ciężkiego sześciokąta nie zapewnia (uwaga własna).
% ŹRÓDŁO: bravyi_high-threshold_2024, streszczenie (12 logical qubits, 288 physical qubits, nearly one million syndrome cycles, 0.1%, nearly 3000; degree-6 graph), sekcja 3 i podpis rys. 1 (kod [[144,12,12]], long-range edges), arXiv:2308.07915v2. Progu nie podano w tekście rozprawy: streszczenie arXiv v2 podaje 0,8%, notatka (streszczenie wersji opublikowanej) 0,7%.
Na platformach o~łączności między dowolnymi kubitami wykonano obliczenia na wielu kubitach logicznych.
Na atomach neutralnych $48$ kubitów logicznych w~blokach $[[8,3,2]]$, zakodowanych w~$128$ kubitach fizycznych, dało z~detekcją błędów wynik lepszy niż kubity fizyczne~\cite{bluvstein_logical_2024}.
W~pracy opublikowanej w~2026~r. kod powierzchniowy o~dystansie~5 miał błąd na rundę $2{,}14(13)$ raza mniejszy niż kod o~dystansie~3~\cite{bluvstein_fault-tolerant_2026}.
% ŹRÓDŁO: bluvstein_logical_2024, streszczenie (48 logical qubits, [[8,3,2]], error detection, outperforming physical qubit fidelities) i część ,,Complex logical circuits using 3D codes'' (128 physical qubits), arXiv:2312.03982v1.
% ŹRÓDŁO: bluvstein_fault-tolerant_2026, streszczenie i podpis rys. 2 (2.14(13)x, four-round circuit, no postselection), arXiv:2506.20661v1 (preprint ma inny tytuł niż wersja opublikowana).
Na 98-kubitowym procesorze jonowym, według preprintów z~2026~r., $48$ kubitów logicznych w~kodzie $[[80,48,4]]$ miało błąd cyklu korekcji mniejszy od fizycznego błędu pamięci~\cite{dasu_computing_2026}, a~w~kodach kolorowych o~dystansie do~7 uzyskano $\Lambda=12^{+3}_{-2}$~\cite{lavasani_quantum_2026}.
% ŹRÓDŁO: dasu_computing_2026, streszczenie (98-qubit Helios, 48--94 logical qubits) i część o cyklu QEC (kod [[80,48,4]]: błąd <= 4e-5 na kubit logiczny wobec 1.1(3)e-4, acceptance rate 0.62(2)), arXiv:2602.22211v1 (preprint).
% ŹRÓDŁO: lavasani_quantum_2026, streszczenie (Lambda = 12 +3/-2, d = 3, 5, 7) i część o eksperymencie pamięci (13, 37 i 73 kubity; dekodowanie offline), arXiv:2610.24077v1 (preprint).
Zastrzeżenia są istotne (zestawienie własne).
Kody o~dystansie 2--4 wymagają odrzucania przebiegów (w~kodzie $[[80,48,4]]$ akceptowano $62\%$), co przenosi koszt z~kubitów na liczbę wykonań (kryterium K3, tabela~\ref{tab:tb03:cost}).
Kod kolorowy o~dystansie~7 zajmował $73$ kubity i~był dekodowany po eksperymencie, a~kody o~wysokim współczynniku kodowania wymagają łączności niedostępnej w~układach nadprzewodzących o~połączeniach między sąsiadami.
% ŹRÓDŁO: łączność: bluvstein_logical_2024, streszczenie (arbitrary connectivity, reconfigurable atom arrays); dasu_computing_2026, streszczenie (high-rate codes require long-range qubit connectivity); zestawienie zastrzeżeń jest własne.

Koszt obliczenia o~użytecznej wielkości pozostaje wysoki: obwód na $100$ kubitach logicznych z~$10^{8}$ bramek~$\mathrm{T}$ wymaga w~kodzie powierzchniowym około $55$~tys. kubitów fizycznych przy błędzie fizycznym $10^{-4}$ i~około $306$~tys. przy $10^{-3}$~\cite{litinski_game_2019}, czyli od około $550$ do około $3000$ na kubit obliczenia (iloraz własny).
% ŹRÓDŁO: litinski_game_2019, sekcja 4 (minimal setup: 164*2*13^2 ~ 55 400 dla p = 1e-4, 210*2*27^2 ~ 306 000 dla p = 1e-3), arXiv:1808.02892v3; ilorazy 554 i 3060 to obliczenie własne.
Oszacowanie kosztu rozkładu liczby 2048-bitowej algorytmem Shora zmalało z~$20$~mln kubitów fizycznych~\cite{gidney_how_2021} do niespełna miliona w~preprincie z~2025~r.~\cite{gidney_how_2025}, przy tych samych założeniach (siatka planarna, błąd bramki $10^{-3}$); do narzutu w~kubitach dochodzi narzut czasowy~\cite{babbush_focus_2021}.
% ŹRÓDŁO: gidney_how_2021, streszczenie (20 million noisy qubits, gate error 1e-3), arXiv:1905.09749v3; gidney_how_2025, streszczenie (less than a million noisy qubits, same assumptions), arXiv:2505.15917v1 (preprint).
% ŹRÓDŁO: babbush_focus_2021, streszczenie (large constant factor overheads associated with error-correction; quadratic speedups will not enable quantum advantage on early generations of such fault-tolerant devices), arXiv:2011.04149v2.

Wniosek dla tej pracy jest zawężony do platformy i~daty.
Na 156-kubitowym procesorze ibm\_fez (podrozdział~\ref{subsec:tb:noise:sources}) zmieściłoby się, przy $2d^2-1$ kubitach na kubit logiczny kodu powierzchniowego~\cite{google_quantum_2025} i~z~pominięciem topologii, najwyżej $9$ kubitów logicznych o~dystansie~3, $3$ o~dystansie~5 albo jeden o~dystansie~7 (wyliczenie własne).
W~eksperymentach z~lat 2022--2023 kody o~dystansie~3 na układach nadprzewodzących miały przy tym błąd od około $3\%$ do $11\%$ na cykl.
% ŹRÓDŁO: wzór 2d^2-1: google_quantum_2025, sekcja I, arXiv:2408.13687v1; 156 kubitów: podrozdział 2.2 rozprawy; 156/17 = 9,2; 156/49 = 3,2; 156/97 = 1,6 (obliczenie własne); 3--11%: krinner2022realizing, google_suppressing_2023, sundaresan_demonstrating_2023 (jak wyżej).
Na urządzeniach nadprzewodzących o~połączeniach między sąsiadami i~około stu kubitach pełna korekcja zużywa zatem, według stanu z~2026~r., cały rejestr na jeden lub kilka kubitów logicznych, a~wykazany zysk dotyczy pamięci, nie obliczeń.
Dla zadań, które jak kodowanie obrazu (rozdziały~\ref{chapter:pdu_images} i~\ref{chapter:learned_encoding}) potrzebują całego rejestru, pozostaje więc niedostępna, choć wyniki z~platform jonowych i~atomowych pokazują, że nie jest to ograniczenie zasadnicze.
